\chapter{Discussion}
\label{ch:discussion}

\section{Interpretation of Results}

\subsection{Co-evolutionary Dynamics}
The positive co-evolutionary gap $\Delta = \Fdef - \Fatt = +0.239 \pm 0.020$
indicates that the defender is winning the arms race at training
completion. However, the gap narrows in episodes 80--120 when the
\tagan{}'s bypass reward plateaus at $\approx 0.055$, suggesting
the attacker has found a locally optimal evasion strategy.
This mirrors real APT behaviour where attackers stabilise on a small
set of effective techniques before pivoting to new vectors.

\subsection{Why RF/DT Outperform \caesar{} on Clean Data}
Random Forest and Decision Tree achieve near-perfect in-distribution
performance (F1 $> 0.993$) on real CICIDS2017 data because the dataset
has clear feature-space separability.
\caesar{}'s lower clean F1 (0.977) reflects the cost of learning a
co-evolutionary policy rather than a static classifier.
The critical advantage of \caesar{} appears under adversarial
conditions (\cref{tab:robustness}): both RF and DT drop from
DR $>$98\% to \textbf{0\%} under \mgae{} perturbation on real data,
while \caesar{}'s adaptive defender maintains its decision quality
because it learned from a dynamically shifting attack distribution.

\subsection{\mgae{} Evasion Analysis}
On real CICIDS2017 data, \mgae{} achieves \textbf{complete evasion}
(0\% detection) of both Random Forest and Decision Tree classifiers,
while FGSM and PGD achieve identical results on these models.
The Threshold IDS remains unaffected (73\% DR) because its z-score
anomaly detection is invariant to the specific perturbation method.
The overall evasion probability of 33.2\% (across all models) is
lower than synthetic-data results, reflecting the harder task of
evading detectors trained on real feature distributions.
The manifold alignment score of 0.001 (vs.\ 0.050 on synthetic data)
indicates that \mgae{} successfully projects attack samples very close
to the benign traffic manifold when trained on real data.

\section{Limitations}

\begin{description}
  \item[Feature-level integration]
    While \caesar{} now operates on real CICIDS2017 data (2.5M flows),
    the co-evolutionary environment maps real flow statistics to simulation
    parameters rather than operating directly on raw PCAP captures.
    Future work should integrate CICFlowMeter-extracted features directly
    into the \tagan{} generator and \adpn{} policy network for
    end-to-end learning on raw network captures.

  \item[Scalability of TAG]
    The current \tagraph{} uses 16 nodes ($8 + 8$).
    Scaling to multi-site enterprise networks with hundreds of attack
    types and countermeasures requires GNN-based reasoning over the
    graph structure rather than simple Markov-chain MLE.

  \item[NumPy vs.\ PyTorch]
    The evolutionary-strategy updates in the NumPy implementation
    are an approximation of true gradient descent. The PyTorch
    implementation (included in \cref{ch:tagan,ch:adpn}) provides
    proper gradient-based training and should be used for
    publication-level experiments.

  \item[Single-site simulation]
    The CyberEnvironment models a single 10-node network.
    Real enterprises are multi-site with inter-site traffic and
    distributed attack campaigns.

  \item[No temporal ordering in MGAE]
    Network flows have temporal structure (IAT distributions, burst
    patterns). The current \mgae{} treats each flow independently.
    A recurrent or attention-based score network would better capture
    temporal correlations in traffic sequences.
\end{description}

\section{Ethical Considerations}

The \caesar{} framework is designed exclusively for defensive
research purposes: training more robust IDS and evaluating the
resilience of security systems before deployment.
All attack generation is performed within a closed simulation
environment with no connection to external networks.
The phishing generation module is designed for detector training
only; generated emails are not distributed.
Release of the source code includes documentation explicitly
prohibiting offensive use.
